\section*{Заключение}
\addcontentsline{toc}{section}{Заключение}

В ходе выполнения курсовой работы было разработано веб-приложение для управления учебными задачами и дедлайнами студентов РТУ МИРЭА с интеграцией официального расписания. Проект решает ключевую проблему --- отсутствие единой среды, которая объединяет расписание пар и персональные задачи студента.

Достигнуты следующие результаты:
\begin{enumerate}
    \item Проведён анализ предметной области и выявлена актуальность автоматизации учёта дедлайнов в контексте учебного расписания.
    \item Спроектирована архитектура на основе Clean Architecture с разделением слоёв домена, бизнес-логики и представления.
    \item Реализован серверный модуль на FastAPI с полнофункциональным REST API и веб-интерфейсом.
    \item Спроектирована база данных пользователей и задач в PostgreSQL с поддержкой миграций Alembic.
    \item Реализована безопасная JWT-аутентификация и хранение паролей с применением современных алгоритмов хэширования.
    \item Создана клиентская часть на Jinja2 с модальными окнами, календарём и автодополнением дисциплин.
    \item Развёрнут отдельный сервис расписания МИРЭА с парсером официального API и хранением данных в MongoDB.
\end{enumerate}

В результате получено готовое к эксплуатации приложение, которое позволяет студентам планировать учебную нагрузку, визуально контролировать дедлайны и связывать их с конкретными учебными занятиями. В дальнейшем возможны расширения: уведомления о сроках, интеграция с мессенджерами и экспорт расписания.
